The Cauchy theory of Four\`es-Bruhat and Choquet-Bruhat--Geroch
\cite{ChoquetBruhat1952,ChoquetBruhatGeroch1969} makes the maximal
globally hyperbolic development the natural spacetime determined by
vacuum initial data. Sbierski \cite{Sbierski2016} gives a modern
construction with explicit initial-surface and embedding hypotheses.
The Kerr solutions \cite{Kerr1963} show why maximality in this class
need not mean inextendibility as a Lorentzian manifold: their analytic
continuation has inner Cauchy horizons, whose global geometry was
analyzed by Carter \cite{Carter1968}. Penrose's cosmic-censorship
programme began with the visibility of singularities in collapse
\cite{Penrose1969}; his later strong formulation linked Cauchy-horizon
instability to global predictability \cite[Section~12.3.2]{Penrose1979}.
The regularity of an extension and the topology used to define
genericity are essential parts of a precise formulation.

The physical blueshift and mass-inflation analyses of Poisson--Israel
and Ori \cite{PoissonIsrael1990,Ori1991,Ori1992} suggested that a
Cauchy horizon can retain a continuous metric while developing a
singularity in its derivatives. Dafermos
\cite{Dafermos2003,Dafermos2005} established rigorous nonlinear
stability and instability results for the spherically symmetric
Einstein--Maxwell--scalar-field equations. In that model,
Luk--Oh \cite{LukOh2019I,LukOh2019II} proved generic \(C^2\)
future inextendibility for admissible complete two-ended asymptotically
flat data, combining exterior lower bounds with interior instability.
Their data space and genericity topologies differ from those used
here. For the vacuum equations without symmetry, Luk
\cite{Luk2018} constructed weak null singularities, and
Dafermos--Luk \cite{DafermosLuk2017} proved continuous extendibility
across a nontrivial Cauchy-horizon portion under quantitative assumptions
on data already inside a black hole. This distinction between metric
continuity and higher regularity is central to the present problem.

Recent results strengthen the vacuum interior instability picture.
Gurriaran \cite{Gurriaran2026} proves curvature blow-up and locally
Lipschitz inextendibility across a horizon portion near timelike
infinity, assuming nonlinear Price-law asymptotics with a nonzero
nonaxisymmetric leading curvature amplitude. Luk--Sbierski
\cite{LukSbierski2026} prove formation of a continuously extendible
but non-Lipschitz weak null singularity from characteristic data with
quantitative decay and a weighted lower bound on a radiative curvature
mode. Sbierski's geometric criteria
\cite{Sbierski2022,Sbierski2026} explain how controlled parallel
transport and robust signed curvature information obstruct Lipschitz
extension. The Kerr interior instability results just cited concern a
prescribed interior or horizon portion under explicit radiation hypotheses. The problem addressed
here is to start with the complete weighted smooth bridge data and
rule out every \(C^2\) future exit generically. The finite perturbations
below provide the needed curvature test violations without imposing a
radiation lower bound on each starting datum.

\paragraph{Exterior evolution.}
The exterior estimates draw on both energy and stationary methods.
The redshift estimate and the physical-space decay method of
Dafermos--Rodnianski
\cite{DafermosRodnianski2009,DafermosRodnianski2010} were followed by
boundedness, decay, and scattering for scalar waves throughout the
subextremal Kerr range
\cite{DafermosRodnianskiShlapentokhRothman2016,DafermosRodnianskiShlapentokhRothman2018}.
For radiative spin-two fields, the Teukolsky estimates of
Dafermos--Holzegel--Rodnianski \cite{DafermosHolzegelRodnianski2019}
and Shlapentokh-Rothman--Teixeira da Costa
\cite{ShlapentokhRothmanTeixeiraDaCosta2023} treat slowly rotating and
full-subextremal Kerr, respectively. At the metric level,
H\"afner--Hintz--Vasy \cite{HafnerHintzVasy2021,HHV2025} establish
linear stability, using stationary spectral analysis and constraint
damping; the full-subextremal mode analysis by
Andersson--H\"afner--Whiting \cite{AnderssonHafnerWhiting2026}
is an important ingredient. These linear results distinguish radiative
decay from the stationary Kerr and gauge modes that must be tracked
in a nonlinear argument.

Nonlinear stability for slowly rotating Kerr was established through
the geometric and wave-equation analysis of Klainerman, Szeftel,
Giorgi, and Shen; see
\cite{KlainermanSzeftel2021,GiorgiKlainermanSzeftel2022} and the
constituent works cited there. Hintz \cite{Hintz2026} establishes
nonlinear exterior stability throughout the subextremal range for
compact perturbations and a broader partially polyhomogeneous
asymptotic class. That class has metric perturbations decaying like
\(r^{-1-\epsilon_0}\), with \(\epsilon_0>0\), and a faster-decaying
remainder; it does not directly include the arbitrary order-\(r^{-1}\)
symbol tails permitted here. We therefore prove the compact and far
estimates for our data class together. Their wave-coordinate
cancellation is related to the weak-null structure identified and
used by Lindblad--Rodnianski
\cite{LindbladRodnianski2003,LindbladRodnianski2010}.
The use of elliptic control at selected frequencies together with
frequency-localized physical-space energy currents also has a close methodological
parallel in the full-subextremal quasilinear-wave analysis of
Dafermos--Holzegel--Rodnianski--Taylor
\cite{DafermosHolzegelRodnianskiTaylor2024}. The tensorial gauge
construction and its estimates are established here for the
Einstein system.

\paragraph{Exact data and finite signals.}
Localized deformation of the constraints was developed by Corvino,
Corvino--Schoen, and Chru\'sciel--Delay
\cite{Corvino2000,CorvinoSchoen2006,ChruscielDelay2003};
Carlotto--Schoen \cite{CarlottoSchoen2016} proved a further conical
localization of asymptotically flat data. Our correction fixes an inner
region while allowing an asymptotic tail, so it can change the charges.
The explicit inverse uses a localized divergence primitive in the
tradition of Bogovskii \cite{Bogovskii1980}, whose support and regularity
properties are treated by Costabel--McIntosh
\cite{CostabelMcIntosh2010}. We prove the required inverse here because its
support, derivative gain, and nonlinear dependence are needed in the
finite experiment. Similarly, complex Gaussian beams and their
superpositions provide the linear antecedents of the oscillatory
construction \cite{Ralston1982,Tanushev2008,LiuRunborgTanushev2013}.
The matching of Einstein wave-gauge data and the bounds on time
intervals growing with the experiment parameter are proved below.
